\section{Evaluation Metrics}
\label{sec:performance_metrics}
In our study, We report Macro Precision, Macro Racall, Macro F1, and Accuracy on the PredEx judgment prediction test dataset and employ a multifaceted approach to evaluate the performance of our models on the PredEx judgment explanation test dataset. Our evaluation metrics encompass both quantitative and qualitative methods, ensuring a thorough assessment of the model's capabilities in both prediction and explanation tasks.

\begin{enumerate}
    \item \textbf{Lexical Based Evaluation:}
    We utilized lexical similarity metrics such as Rouge scores (Rouge-1, Rouge-2, and Rouge-L) ~\cite{lin-2004-rouge}, BLEU ~\cite{papineni-etal-2002-bleu}, and METEOR ~\cite{banerjee-lavie-2005-meteor}. These metrics assess the similarity between the generated explanations and the reference texts based on word overlap and order, providing an insight into the lexical accuracy of the model outputs.

    \item \textbf{Semantic Similarity Based Method:}
    To capture the semantic essence of the generated explanation, we employed BERTScore ~\cite{BERTScore}, which measures the semantic similarity between the generated and ground truth explanations. Additionally, we used BLANC ~\cite{blanc} to estimate the quality of generated explanations in the absence of a gold standard, offering a perspective on the model's ability to generate semantically rich and contextually relevant text.

    \item \textbf{Expert Evaluation:}
    Human evaluation played a crucial role in our assessment. Legal experts reviewed the explanations generated by the models and rated them on a 1–5 Likert scale based on their accuracy, relevance, and completeness. The criteria for the rating scale were as follows:
    \begin{enumerate}[label=\arabic*.]
        \item The explanation is entirely incorrect or fails to provide any relevant information.
        \item The model's response is irrelevant or shows misunderstanding of the case judgment.
        \item The explanation is partially accurate but misses critical details.
        \item The response is comparable and relevant to the ground truth.
        \item The explanation is completely accurate, relevant, and potentially superior to the expert's explanation.
    \end{enumerate}
    
    % \item \textbf{Statistical Significance:}
    % We conducted statistical analysis to determine the significance of the differences observed in the performance of various models. This involved comparing BERT similarity scores across different experimental settings of the LLMs. A p-value of 0.05 or lower {\color{red}KRIPA: more details and citation needed. t-test? two-sided? significance level?} was considered indicative of statistical significance. Such a finding implies that the observed differences in performance are unlikely to be attributable to random chance.
\end{enumerate}